\documentclass[12pt]{article}
\usepackage{amsmath, amssymb}
\usepackage{fancyhdr}
\usepackage{fullpage}
\usepackage{listings}
\usepackage{enumitem}
\usepackage{graphicx}
\usepackage{courier}
\usepackage{float}
\usepackage{hyperref}

\pagestyle{fancy}
\fancyhf{}
\rhead{CPSC 250 - Final Exam}
\lhead{Programming for Data Manipulation}
\rfoot{Page \thepage}

\title{CPSC 250 - Programming for Data Manipulation\\Final Exam}
\date{}
\begin{document}
\maketitle

\section*{Instructions}
You have 150 minutes. This is a closed-book, in-class exam. No access to the internet or external software is allowed. Show all work clearly and legibly.

\section*{Part I: Multiple Choice (30 points, 2 points each)}
Circle the best answer.

\begin{enumerate}[label=\arabic*.]

\item What is the output of the following code?  
\texttt{a = [1, 2, 3] \\
b = a \\
b[0] = 99 \\
print(a)}  
\begin{enumerate}[label=\Alph*.]
\item \texttt{[1, 2, 3]}  
\item \texttt{[99, 2, 3]}  
\item \texttt{[1, 99, 3]}  
\item Error
\end{enumerate}

\item What does \texttt{df.dropna()} do in pandas?  
\begin{enumerate}[label=\Alph*.]
\item Fills in missing values  
\item Drops rows or columns with missing values  
\item Drops duplicate rows  
\item Drops specified columns
\end{enumerate}

\item What is the purpose of the \texttt{\_\_str\_\_} method in a class?  
\begin{enumerate}[label=\Alph*.]
\item Initialize the object  
\item Represent the object as a string  
\item Destroy the object  
\item Compare two objects
\end{enumerate}

\item What is the primary benefit of inheritance in object-oriented programming?  
\begin{enumerate}[label=\Alph*.]
\item Improved speed  
\item Data hiding  
\item Code reuse  
\item Loop simplification
\end{enumerate}

\item In a recursive function, what is the base case?  
\begin{enumerate}[label=\Alph*.]
\item The case where an exception is raised  
\item The function that calls itself  
\item The condition where recursion stops  
\item The loop inside the function
\end{enumerate}

\item Which of the following is true about lists and tuples?  
\begin{enumerate}[label=\Alph*.]
\item Lists are immutable, tuples are mutable  
\item Both are mutable  
\item Both are immutable  
\item Lists are mutable, tuples are immutable
\end{enumerate}

\item What does the \texttt{curve\_fit} function from \texttt{scipy.optimize} do?  
\begin{enumerate}[label=\Alph*.]
\item Plots data using a curve  
\item Fits data to a specified function  
\item Fits a line to data using OLS  
\item Returns the correlation coefficient
\end{enumerate}

\item What is the time complexity of searching in a balanced binary search tree?  
\begin{enumerate}[label=\Alph*.]
\item \texttt{O(1)}  
\item \texttt{O(n)}  
\item \texttt{O(log n)}  
\item \texttt{O(n log n)}
\end{enumerate}

\item Which of the following statements creates a dictionary?  
\begin{enumerate}[label=\Alph*.]
\item \texttt{d = \{\}}  
\item \texttt{d = []}  
\item \texttt{d = ()}  
\item \texttt{d = set()}
\end{enumerate}

\item What happens when you override the \texttt{\_\_eq\_\_} method in a class?  
\begin{enumerate}[label=\Alph*.]
\item You redefine how two instances are compared with \texttt{==}  
\item You make your class hashable  
\item You initialize the object  
\item You compare types of objects
\end{enumerate}

\item What does the pandas method \texttt{read\_csv()} return?  
\begin{enumerate}[label=\Alph*.]
\item A list of strings  
\item A NumPy array  
\item A Series  
\item A DataFrame
\end{enumerate}

\item What does operator overloading allow in Python?  
\begin{enumerate}[label=\Alph*.]
\item Redefining loops  
\item Creating multiple classes in one file  
\item Using operators like + with user-defined types  
\item Automatically generating documentation
\end{enumerate}

\item In pandas, what does \texttt{df['column'].mean()} compute?  
\begin{enumerate}[label=\Alph*.]
\item The median  
\item The total  
\item The standard deviation  
\item The average
\end{enumerate}

\item What is polymorphism in OOP?  
\begin{enumerate}[label=\Alph*.]
\item The ability of different classes to share the same attributes  
\item The ability of objects to take on multiple forms via shared methods  
\item Hiding data inside an object  
\item Changing variable types at runtime
\end{enumerate}

\item What is the correct way to inherit from a class in Python?  
\begin{enumerate}[label=\Alph*.]
\item \texttt{class A inherits B:}  
\item \texttt{class A: inherits B}  
\item \texttt{class A(B):}  
\item \texttt{class A<B>:}
\end{enumerate}

\end{enumerate}

\newpage
\section*{Part II: Error Identification (15 points)}

The following code is meant to reverse a list recursively. Identify and correct all errors.

\begin{lstlisting}[language=Python]
def reverse_list(lst):
    if len(lst) = 0:
        return []
    else:
        return reverse_list(lst[1:]) + lst[0]
\end{lstlisting}

\vspace{3in}

\newpage
\section*{Part III: Code Writing (30 points)}

\textbf{Q1. (6 points)}  
Write a function \texttt{is\_palindrome(s)} that returns \texttt{True} if the string \texttt{s} is a palindrome, and \texttt{False} otherwise.

\vspace{2in}

\textbf{Q2. (8 points)}  
Write a class \texttt{Student} with fields for \texttt{name} and \texttt{gpa}, and a \texttt{\_\_str\_\_} method that prints the student’s name and GPA. Also include an \texttt{\_\_lt\_\_} method that compares students by GPA.

\vspace{2.5in}

\textbf{Q3. (8 points)}  
Write a function \texttt{load\_and\_filter(filename)} that reads a CSV file into a pandas DataFrame, and returns only the rows where the value in the column \texttt{"score"} is greater than 80.

\vspace{2.5in}

\textbf{Q4. (8 points)}  
Write a base class \texttt{Animal} with a method \texttt{make\_sound()} that returns the string \texttt{"..."}. Create two derived classes \texttt{Dog} and \texttt{Cat} that override \texttt{make\_sound()} to return \texttt{"woof"} and \texttt{"meow"}, respectively. Show how polymorphism allows you to loop over a list of animals and call \texttt{make\_sound()} on each.

\vspace{3in}

\newpage
\section*{Part IV: Code Comprehension and Commenting (25 points)}

This function is designed to perform curve fitting using \texttt{curve\_fit}. Comment each line and explain what the code is doing overall.

\begin{lstlisting}[language=Python]
import numpy as np
from scipy.optimize import curve_fit

def model(x, a, b):
    return a * np.exp(b * x)

def fit_data(x, y):
    popt, pcov = curve_fit(model, x, y)
    print("a =", popt[0])
    print("b =", popt[1])
\end{lstlisting}

\vspace{3in}

\textbf{What do \texttt{popt} and \texttt{pcov} represent in this context?}

\vspace{2in}

\end{document}
